\documentclass[11pt]{article}

\usepackage{amsmath,amssymb}
\usepackage{xurl}
\usepackage[hidelinks]{hyperref}
\setlength{\emergencystretch}{2em}

% Shared commands for notes in docs/paper-gaps/.
%
% Two halves.  The link commands below are project identity: OWNER/REPO is the
% placeholder that scripts/bootstrap_project.py replaces with this project's
% GitHub slug and Pages site.  Everything after them is NOTATION, and notation
% belongs to the source paper: the definitions here are an example set, and a
% project replaces them with the macros of its own paper's style file.  The one
% rule that never changes: a command defined here must mean exactly what the
% source paper's macro means, or a note silently says something else.

\newcommand{\Lean}{\textsc{Lean}}
\newcommand{\leanid}[1]{\textsf{#1}}
\newcommand{\ghissue}[1]{\href{https://github.com/OWNER/REPO/issues/#1}{GitHub issue \##1}}
\newcommand{\ghpr}[1]{\href{https://github.com/OWNER/REPO/pull/#1}{PR \##1}}
% Local issue tracker (issues/NNNN-*.md in this repository); no remote URL.
\newcommand{\localissue}[1]{issue \##1 (\path{issues/})}
\newcommand{\doclink}[1]{\href{https://OWNER.github.io/REPO/docs/#1.html}{\path{#1}}}

% --- Notation: replace with the source paper's own macros --------------------
% \F, \N, \C, \R, \tr, \ot, \eps, \E, \bx, \bu, \bv, \by

\newcommand{\F}{\mathbb{F}}
\newcommand{\N}{\mathbb{N}}
\newcommand{\C}{\mathbb{C}}
\newcommand{\R}{\mathbb{R}}
\newcommand{\tr}{\mathrm{tr}}
\newcommand{\eps}{\epsilon}
\newcommand{\ot}{\otimes}
\newcommand{\E}{\mathop{\bf E\/}}

% bold random variables
\newcommand{\bu}{\boldsymbol{u}}
\newcommand{\bv}{\boldsymbol{v}}
\newcommand{\bx}{{\boldsymbol{x}}}
\newcommand{\by}{\boldsymbol{y}}

% T1 so literal < and > in \texttt placeholders typeset as themselves
\usepackage[T1]{fontenc}

% bra-ket
\usepackage{braket}

% --- Example structured spaces ---
\newcommand{\polyfunc}[3]{\mathcal{P}(#1, #2, #3)}
\newcommand{\polymeas}[3]{\mathrm{PolyMeas}(#1, #2, #3)}
\newcommand{\polysub}[3]{\mathrm{PolySub}(#1, #2, #3)}

% Approximation relations (paper uses \approx_\eps and \simeq_\zeta)
\newcommand{\approxeps}{\approx_{\varepsilon}}
\newcommand{\simgeq}{\simeq_{\zeta}}

\renewcommand{\ghissue}[1]{%
  \href{https://github.com/Dengnifer/MIPStarRE-A/issues/#1}{GitHub issue \##1}}

\title{Triangle Estimates and Isometry Transfer in Pauli Extraction}
\author{}
\date{2026-09-15}

\begin{document}
\maketitle

\section{At a glance}

\begin{itemize}
  \item \textbf{Difficulty.} Medium.  The extraction proof contains two
  incorrect numerical steps and later suppresses the range projection when
  passing from unitary conjugation to isometry conjugation.
  \item \textbf{Estimated weight.} A corrected positive-contraction estimate,
  a small-error case split for normalization, and one range-projection transfer
  lemma for a projective measurement family.
  \item \textbf{Mathlib/project split.} Positivity, contraction inequalities,
  isometry range projections, and norm estimates belong to finite-dimensional
  Hilbert-space algebra.  The state-dependent family distance and register
  placements are project-local.
  \item \textbf{Key mathematical inputs.} The inequality
  $H^2\geq2H-I$ for a positive contraction, the triangle inequality, and
  $Q=\phi\phi^\dagger$ for an isometry.  The repair is tracked in
  \ghissue{19}.
\end{itemize}

\section{Key theorem forms}

Let $H_X$ and $H_Z$ be the positive contractions of the extraction proof, and
let $\vartheta$ be a unit vector such that
\[
  \langle\vartheta,H_W\vartheta\rangle\geq1-\delta/2
  \qquad (W\in\{X,Z\}).
\]
Since $(I-H_W)^2\geq0$,
\[
  H_W^2\geq2H_W-I,
  \qquad
  \langle\vartheta,H_W^2\vartheta\rangle\geq1-\delta.
\]
Also $\lVert(I-H_W)\vartheta\rVert^2\leq\delta$, whence
\[
  \lVert H_X\vartheta-H_Z\vartheta\rVert^2\leq4\delta.
\]
Together with $H_XH_Z$ equal to the EPR projection, these estimates give the
required overlap bound, with constants absorbed into the stated
$O(\sqrt\delta)$ error.

For the final transfer, let $\phi\colon\mathcal H\to\mathcal K$ be an isometry
and set $Q=\phi\phi^\dagger$.  If $N_h$ is the unitary-conjugated extension of
a measurement operator $M_h$, then the exact identity is
\[
  \phi M_h\phi^\dagger=N_hQ=QN_h,
\]
not merely $N_h$.  Thus comparison of $N_h$ with an ideal projector must be
combined with control of $(I-Q)\Delta$ on the ideal comparison state
$\Delta$.

\section{Numerical defects in the source proof}

Lines~1715--1742 of
\path{references/qpbt-paper/14_analysis_of_the_pauli_basis_test.tex}
construct $H_X,H_Z$ and identify $H_XH_Z$ with the EPR projection.  The norm
estimate at lines~1743--1750 gives
$\lVert(I-H_W)\vartheta\rVert\leq\sqrt\delta$.  At lines~1753--1757 the
source substitutes $\delta$ inside the reverse triangle inequality
\cite[lines~1715--1783]{Ji2020MIPStar}.  That replacement is not valid for
small $\delta$.  The desired lower bound on
$\langle H_W^2\rangle$ follows directly from $H_W^2\geq2H_W-I$, as above.

The triangle inequality yields
\[
  \lVert H_X\vartheta-H_Z\vartheta\rVert\leq2\sqrt\delta,
\]
so its square is at most $4\delta$.  Line~1761 instead bounds that square by
$2\sqrt\delta$ without a stated restriction.  This inequality follows when
$\delta\leq1/4$; for larger error the theorem is made trivial by enlarging the
universal constant.  The same small-error split ensures that the projected
vector defined at lines~1769--1773 is nonzero before it is normalized.  These
are local numerical repairs and do not weaken the extraction conclusion.

\section{The range-projection obligation}

Lines~1827--1858 prove state-dependent closeness for the concrete operators
obtained by conjugating the original measurement with the swap unitaries.
Lines~1864--1875 then define isometries by adjoining EPR pairs and applying
those unitaries, and immediately state the corresponding closeness for
$\phi M_h\phi^\dagger$.  The omitted range projection makes this implication
non-definitional \cite[lines~1827--1875]{Ji2020MIPStar}.

Let $\Psi$ be the original state and suppose $\phi\Psi$ is close to the ideal
state $\Delta$.  Since $Q\phi\Psi=\phi\Psi$ and $Q$ is a contraction, the
triangle inequality bounds $(I-Q)\Delta$ by the state-extraction error.  Using
$\phi M_h\phi^\dagger=N_hQ$, split
\[
  (N_hQ-T_h)\Delta
  =(N_h-T_h)\Delta-N_h(I-Q)\Delta,
\]
where $T_h$ is the ideal Pauli projector.  Summing squared norms over $h$
uses projectivity and orthogonality to keep the second term under control.
The same argument is required for each player and each Pauli basis.

This transfer is an open proof obligation for the next extraction stage.  It
must be a theorem derived from the concrete swap unitaries, transformed-state
estimate, and operator comparisons.  It must not be stored as a bridge field,
residual assumption, or hypothesis of the source-labelled soundness theorem.

\section{Blueprint and formalization status}

The proof of Lemma~\path{lem:qld-unitary} and the soundness proof in
\path{blueprint/src/chapter/ch16_qpbt_extraction.tex} use the corrected
positive-contraction estimate, make the small-error case distinction explicit,
and isolate the range-projection calculation.  The chapter therefore preserves
the source's state and observable conclusions while exposing the omitted
argument.

The structure \path{ExtractionWitness} in
\path{MIPStarRE/QPBT/Extraction/Unitary.lean} contains only the two-sided
unitarity equations for the concrete swap maps, a normalized auxiliary state,
the transformed-state estimate, and the two-player Pauli-projector comparison.
The theorem \path{exists_extractionWitness_ofGlobalPairWitness} concludes the
existence of this data at the explicit extraction scale, but it takes the
\path{GlobalPairWitness} of Lemma~\path{lem:qld-4-7} as an explicit hypothesis.
It therefore remains a conditional helper rather than a formalization of the
source lemma.  The blueprint records it on the separate formalization-support
node \path{lem:qld-unitary-given-global-pair}, whose statement displays the
global polynomial-pair measurement as supplied data.  The helper is proved at
this conditional statement and has only the standard axioms \path{propext},
\path{Classical.choice}, and \path{Quot.sound}.  Apart from the supplied witness,
its conclusion matches the Stage~4.2 extraction lemma, including the corrected
numerical estimates.

The preceding global-measurement construction is now complete.
The theorem \path{exists_globalPairWitness} proves the full-domain construction
of Lemma~\path{lem:qld-4-7} for every strategy error, including zero.  The
source-facing theorem \path{exists_pulled_apart_consistency} chooses one such
witness, applies the two point-consistency estimates and the observable
self-consistency estimate to it, and replaces their three universal constants
by their maximum.  Thus Lemma~\path{lem:qld-construct-the-paulis} is proved for
the stated domain $0\leq\epsilon\leq1$, including the zero-error boundary.

The source-facing theorem \path{exists_extractionWitness} now composes
\path{exists_globalPairWitness} with
\path{exists_extractionWitness_ofGlobalPairWitness}: the universal error
constants for the global measurement and extraction are chosen before the
parameters and strategy. For every admissible parameter tuple and projective
strategy with $0\leq\epsilon\leq1$, it constructs the global pair of
projective measurements internally and obtains a normalized auxiliary state
and two concrete swap unitaries satisfying the state and both players'
Pauli-family distances at the nested construction/extraction scale. This is
the existence claim of Lemma~\path{lem:qld-unitary}; its proof does not
require a supplied global measurement. The subsequent passage from the
concrete swap operators to the isometries in the soundness theorem remains
unformalized and still requires the range-projection argument above.

Issue~\#19 owns this concrete Stage~4.2 interface and discrepancy record.
\ghissue{123} records the source-facing composition work: both the pulled-apart
consistency and the source-unitary existence compositions are proved. Other
issue~\#123 targets remain unfinished; the earlier attempt history is unchanged.
The direct-carrier proof of the global constructor does not discharge the
seed-indexed correspondence and direct-law restrictions documented in
\path{docs/paper-gaps/qpbt_ld-dimension-divisibility.tex}.
\ghissue{47} tracks the separate range-projection transfer needed
by \path{pauli_soundness}; that theorem must be derived from the concrete data
rather than stored in the witness.

\section{Conclusion}

The printed triangle calculation has repairable numerical errors, and the
last passage from unitaries to isometries omits a genuine range-projection
argument.  Neither defect is a counterexample to Pauli soundness.  The former
is corrected by positive-contraction inequalities and a small-error split;
the latter remains a named downstream proof obligation whose hypotheses are
already supplied by the extraction lemma.  The global polynomial-pair witness
and the three pulled-apart consistency estimates are now composed in a proved
source-facing theorem, including at zero error. The conditional extraction
theorem remains useful for a supplied global measurement, and its composition
with the proved global constructor establishes the source-facing existence
claim of Lemma~\path{lem:qld-unitary} for projective strategies. The final
range-projection transfer remains separate and open.

\bibliographystyle{alpha}
\bibliography{references}

\end{document}
